\begin{chaptersummary}
  \item Hot Chocolate derives the schema from your C\# types at compile time.
        \code{[QueryType]} feeds a source generator, which writes a
        registration method named after your assembly, so the line you call is
        \code{AddSessionsTypes} here and something else in your project.
  \item Leaving that attribute off fails in two opposite ways. On the last
        remaining root class the host refuses to start and names the problem;
        on any other class the build passes, the service starts, and the
        fields are absent.
  \item Your nullable reference type annotations are the schema's non-null
        annotations. \code{string} becomes \code{String!} and \code{string?}
        becomes \code{String}, and a client's generated code will believe both.
  \item Every public member is published unless you mark it otherwise, so a
        foreign key added for the compiler becomes an API commitment.
        \code{[GraphQLIgnore]} is the only thing between \code{SpeakerId} and
        every client you will ever have.
  \item The loop is export the schema, read the diff, send the request. The
        schema is validated against before any resolver runs, which is why a
        misspelled field costs a round trip rather than an afternoon.
\end{chaptersummary}
